\documentclass[10pt,a4paper]{article}
\usepackage[utf8]{inputenc}
\usepackage[spanish]{babel}
\usepackage[left=1cm, right=1cm, top=2cm, bottom=1cm, headheight=25pt]{geometry}
\usepackage{multicol}
\usepackage{booktabs}
\usepackage{tabularx}
\usepackage[table]{xcolor}
\usepackage{amsmath}
\usepackage{titlesec}
\usepackage{fancyhdr}

\pagestyle{fancy}
\fancyhf{}
\fancyhead[L]{UTN-FRC \\ Técnicas Digitales II}
\fancyhead[C]{\textbf{Hoja de Referencia BluePill} \\ \textbf{STM32F103C8T6}}
\fancyhead[R]{Año 2026 \\ Franco Palombo \quad Pág. \thepage}
\renewcommand{\headrulewidth}{0.4pt}
\renewcommand{\footrulewidth}{0pt}

\newcolumntype{L}{>{\raggedright\arraybackslash}X}

% Colores de la tabla similar al PDF
\definecolor{tableheadcolor}{HTML}{D9E1F2}
\definecolor{tablelightcolor}{HTML}{F2F2F2}

% Formato de secciones
\titleformat{\section}{\large\bfseries\color{blue!80!black}}{}{0pt}{}[\vspace{-2pt}\titlerule]
\titlespacing*{\section}{0pt}{10pt}{5pt}

\setlength{\parindent}{0pt}

% Comando para fila sombreada
\newcommand{\rowgray}{\rowcolor{tablelightcolor}}
\newcommand{\rowhead}{\rowcolor{tableheadcolor}}

\begin{document}

\small
\begin{multicols*}{2}

\section*{GPIO}
\renewcommand{\arraystretch}{1.2}
\begin{tabularx}{\columnwidth}{l L}
\rowhead \textbf{Elemento} & \textbf{Definiciones CMSIS} \\
\textbf{Clock} & \texttt{RCC->APB2ENR} \\
\rowgray \textbf{Puertos} & \texttt{GPIOA}, \texttt{GPIOB}, \texttt{GPIOC} \\
\textbf{Configuración} & \texttt{CRL}, \texttt{CRH} \\
\rowgray \textbf{Entrada/salida} & \texttt{IDR}, \texttt{ODR}, \texttt{BSRR}, \texttt{BRR} \\
\textbf{Acceso} & \texttt{GPIOx->CRL}, \texttt{GPIOx->CRH} \\
 & \texttt{GPIOx->IDR}, \texttt{GPIOx->ODR} \\
\rowgray \textbf{Máscaras} & \texttt{RCC\_APB2ENR\_IOPxEN} \\
\rowgray & \texttt{RCC\_APB2ENR\_AFIOEN} \\
\rowgray & \texttt{GPIO\_IDR\_IDRp} \\
\rowgray & \texttt{GPIO\_ODR\_ODRp} \\
\end{tabularx}

\vspace{1mm}
\textbf{Configuración de Puertos} \\
\begin{tabularx}{\columnwidth}{|c|c|L|}
\hline
\rowhead \textbf{CNF} & \textbf{MODE} & \textbf{Configuración (Modo)} \\
\hline
00 & 00 & Entrada analógica \\
01 & 00 & Entrada flotante (Reset) \\
10 & 00 & Entrada pull-up / pull-down \\
11 & 00 & \textit{Reservado} \\
\hline
00 & $>00$ & Salida G.P. Push-pull \\
01 & $>00$ & Salida G.P. Open-drain \\
10 & $>00$ & Salida Alt. Fun. (AF) Push-pull \\
11 & $>00$ & Salida Alt. Fun. (AF) Open-drain \\
\hline
\multicolumn{3}{|l|}{\cellcolor{tableheadcolor}\textbf{Opciones de MODE (Velocidad)}} \\
\hline
-- & 00 & Modo Entrada (Reset) \\
-- & 01 & Salida, velocidad máx 10 MHz \\
-- & 10 & Salida, velocidad máx 2 MHz \\
-- & 11 & Salida, velocidad máx 50 MHz \\
\hline
\end{tabularx}
\footnotesize
\textbf{Nota:} \texttt{CRy} representa \texttt{CRL} (pines 0-7) o \texttt{CRH} (8-15). \texttt{p} es el número del pin. \texttt{x} es el puerto (A, B, C). \texttt{>00} significa 01, 10 o 11.

\vspace{1mm}
\textbf{Mapeo de Pines de Periféricos} \\
\begin{tabularx}{\columnwidth}{l L}
\rowhead \textbf{Periférico} & \textbf{Pin y Modo de Configuración} \\
\textbf{USART1 TX} & \texttt{PA9} -- Salida Alt. Fun. Push-Pull \\
\rowgray \textbf{USART1 RX} & \texttt{PA10} -- Entrada Flotante / Pull-Up \\
\textbf{TIM2 (CH1..4)} & \texttt{PA0}..\texttt{PA3} -- Salida Alt. Fun. Push-Pull \\
\rowgray \textbf{ADC1 (IN0..IN7)} & \texttt{PA0}..\texttt{PA7} -- Entrada Analógica \\
\textbf{LED Onboard} & \texttt{PC13} -- Salida G.P. Push-Pull \\
\end{tabularx}

\section*{EXTI}
\begin{tabularx}{\columnwidth}{l L}
\rowhead \textbf{Parámetro} & \textbf{Valor / Registro / Máscara} \\
\rowgray \textbf{Selección} & \texttt{AFIO->EXTICR[0] a [3]} \\
\textbf{Registros} & \texttt{EXTI->IMR}, \texttt{EXTI->RTSR} \\
& \texttt{EXTI->FTSR}, \texttt{EXTI->PR} \\
\rowgray \textbf{IRQ Indiv.} & \texttt{EXTI0\_IRQn} a \texttt{EXTI4\_IRQn} \\
\textbf{IRQ Agrupadas} & \texttt{EXTI9\_5\_IRQn}, \texttt{EXTI15\_10\_IRQn} \\
\rowgray \textbf{ISR usuales} & \texttt{EXTI0\_IRQHandler(void)} \\
\rowgray & a \texttt{EXTI15\_10\_IRQHandler(void)} \\
\textbf{Máscaras} & \texttt{EXTI\_IMR\_MRi}, \texttt{EXTI\_RTSR\_TRi} \\
& \texttt{EXTI\_FTSR\_TRi}, \texttt{EXTI\_PR\_PRi} \\
\end{tabularx}
\footnotesize
\textbf{Nota:} \texttt{i} representa la línea EXTI (0 a 15).

\section*{NVIC}
\begin{tabularx}{\columnwidth}{l L}
\rowhead \textbf{Parámetro} & \textbf{Función CMSIS} \\
\rowgray \textbf{Funciones} & \texttt{NVIC\_EnableIRQ(IRQn\_Type)} \\
\rowgray & \texttt{NVIC\_SetPriority(IRQn\_Type, uint32\_t)} \\
\end{tabularx}


\section*{ADC1}
\begin{tabularx}{\columnwidth}{l L}
\rowhead \textbf{Elemento} & \textbf{Definiciones CMSIS} \\
\textbf{Clock} & \texttt{RCC->APB2ENR} \\
& \texttt{RCC\_APB2ENR\_ADC1EN} \\
\rowgray \textbf{Periférico} & \texttt{ADC1} \\
\textbf{Control} & \texttt{ADC1->CR1}, \texttt{ADC1->CR2} \\
\rowgray \textbf{Muestreo} & \texttt{ADC1->SMPR1}, \texttt{ADC1->SMPR2} \\
\textbf{Secuencia} & \texttt{ADC1->SQR1}, \texttt{ADC1->SQR2} \\
& \texttt{ADC1->SQR3} \\
\rowgray \textbf{Estado/dato} & \texttt{ADC1->SR}, \texttt{ADC1->DR} \\
\textbf{Máscaras CR2} & \texttt{ADC\_CR2\_ADON} \\
& \texttt{ADC\_CR2\_SWSTART} \\
& \texttt{ADC\_CR2\_CONT} \\
& \texttt{ADC\_CR2\_TSVREFE} \\
& \texttt{ADC\_CR2\_EXTSEL} \\
& \texttt{ADC\_CR2\_EXTTRIG} \\
\rowgray \textbf{Otras} & \texttt{ADC\_SR\_EOC}, \texttt{ADC\_SMPR2\_SMP2} \\
\end{tabularx}

\section*{UART -- USART1}
\begin{tabularx}{\columnwidth}{l L}
\rowhead \textbf{Elemento} & \textbf{Definiciones CMSIS} \\
\textbf{Clock} & \texttt{RCC\_APB2ENR\_USART1EN} \\
\rowgray \textbf{Periférico} & \texttt{USART1} \\
\textbf{Registros} & \texttt{USART1->SR}, \texttt{USART1->DR} \\
& \texttt{USART1->BRR} \\
& \texttt{USART1->CR1}, \texttt{USART1->CR2} \\
\rowgray \textbf{Habilitación} & \texttt{USART\_CR1\_UE}, \texttt{USART\_CR1\_TE} \\
\rowgray & \texttt{USART\_CR1\_RE} \\
\textbf{Recepción} & \texttt{USART\_CR1\_RXNEIE} \\
& \texttt{USART1\_IRQn} \\
& \texttt{USART1\_IRQHandler(void)} \\
\rowgray \textbf{Flags} & \texttt{USART\_SR\_TXE}, \texttt{USART\_SR\_RXNE} \\
\rowgray & \texttt{USART\_SR\_TC} \\
\end{tabularx}
\begin{equation*}
    \text{Baud} = \frac{f_{\text{CK}}}{16 \times \text{USARTDIV}}
\end{equation*}

\section*{PWM -- TIM2, canal c}
\begin{tabularx}{\columnwidth}{l L}
\rowhead \textbf{Elemento} & \textbf{Definiciones CMSIS} \\
\textbf{Clock} & \texttt{RCC\_APB1ENR\_TIM2EN} \\
\rowgray \textbf{Periférico} & \texttt{TIM2} \\
\textbf{Base tiempo} & \texttt{TIM2->PSC}, \texttt{TIM2->ARR} \\
\rowgray \textbf{Comparación} & \texttt{TIM2->CCRc} \\
\textbf{Salida} & \texttt{TIM2->CCMRz}, \texttt{TIM2->CCER} \\
\rowgray \textbf{Eventos} & \texttt{TIM2->CR1}, \texttt{TIM2->EGR} \\
\textbf{Máscaras} & \texttt{TIM\_CR1\_CEN} \\
& \texttt{TIM\_CCER\_CCxE} \\
& \texttt{TIM\_CCER\_CCxP} \\
& \texttt{TIM\_CCMRz\_OCxPE} \\
& \texttt{TIM\_CCMRz\_OCxM} \\
& \texttt{TIM\_EGR\_UG} \\
\end{tabularx}
\begin{equation*}
    f_{\text{PWM}} = \frac{f_{\text{CK\_PSC}}}{(\text{PSC}+1)(\text{ARR}+1)} \quad \text{Duty} = \frac{\text{CCR}x}{\text{ARR}+1} \times 100\%
\end{equation*}
\footnotesize
\textbf{Nota:} \texttt{c} o \texttt{x} es el nro de canal (1 a 4). \texttt{z} es 1 para canales 1 y 2; es 2 para canales 3 y 4.

\section*{SysTick}
\begin{tabularx}{\columnwidth}{l L}
\rowhead \textbf{Elemento} & \textbf{Definiciones CMSIS} \\
\textbf{Registros} & \texttt{SysTick->CTRL}, \texttt{SysTick->LOAD} \\
& \texttt{SysTick->VAL} \\
\rowgray \textbf{Máscaras} & \texttt{SysTick\_CTRL\_ENABLE\_Msk} \\
\rowgray & \texttt{SysTick\_CTRL\_TICKINT\_Msk} \\
\rowgray & \texttt{SysTick\_CTRL\_CLKSOURCE\_Msk} \\
\rowgray & \texttt{SysTick\_CTRL\_COUNTFLAG\_Msk} \\
\textbf{Handlers} & \texttt{SysTick\_Config(uint32\_t ticks)} \\
& \texttt{SysTick\_Handler(void)} \\
\end{tabularx}

\end{multicols*}

\newpage
\small
\begin{multicols*}{2}

% Definición de estilo para listas compactas
\let\olditemize\itemize
\renewcommand{\itemize}{\olditemize\setlength{\itemsep}{1pt}\setlength{\parskip}{0pt}\setlength{\parsep}{0pt}}

\section*{RCC (\texttt{0x4002 1000})}
\begin{itemize}
    \item \textbf{APB2ENR} (Offset \texttt{0x18}) -- \textit{APB2 Peripheral Clock Enable Register}
    \begin{itemize}
        \item \texttt{[14]} \texttt{USART1EN}: Reloj USART1.
        \item \texttt{[9]} \texttt{ADC1EN}: Reloj ADC1.
        \item \texttt{[4]} \texttt{IOPCEN}: Reloj Puerto C.
        \item \texttt{[3]} \texttt{IOPBEN}: Reloj Puerto B.
        \item \texttt{[2]} \texttt{IOPAEN}: Reloj Puerto A.
        \item \texttt{[0]} \texttt{AFIOEN}: Reloj Funciones Alternativas.
    \end{itemize}
    \item \textbf{APB1ENR} (Offset \texttt{0x1C}) -- \textit{APB1 Peripheral Clock Enable Register}
    \begin{itemize}
        \item \texttt{[0]} \texttt{TIM2EN}: Reloj Timer 2.
    \end{itemize}
\end{itemize}

\section*{GPIO (\texttt{0x4001} A=\texttt{0800}, B=\texttt{0C00}, C=\texttt{1000})}
\textit{Relativas a \texttt{0x4001 0000}}
\begin{itemize}
    \item \textbf{CRL / CRH} (Offsets \texttt{0x00} / \texttt{0x04}) -- \textit{Port Config Reg Low / High}
    \begin{itemize}
        \item \texttt{[4p+3 : 4p+2]} \texttt{CNFp}: Config. del pin.
        \item \texttt{[4p+1 : 4p]} \texttt{MODEp}: Modo/Velocidad del pin.
    \end{itemize}
    \item \textbf{IDR / ODR} (Offsets \texttt{0x08} / \texttt{0x0C}) -- \textit{Input / Output Data Register}
    \begin{itemize}
        \item \texttt{[15:0]}: Estado de entrada (\texttt{IDR}) o salida (\texttt{ODR}).
    \end{itemize}
    \item \textbf{BSRR} (Offset \texttt{0x10}) -- \textit{Bit Set/Reset Register}
    \begin{itemize}
        \item \texttt{[31:16]} \texttt{BRy}: Reset pin $y$ (escribir 1 pone pin en 0).
        \item \texttt{[15:0]} \texttt{BSy}: Set pin $y$ (escribir 1 pone pin en 1).
    \end{itemize}
    \item \textbf{BRR} (Offset \texttt{0x14}) -- \textit{Bit Reset Register}
\end{itemize}

\section*{AFIO (\texttt{0x4001 0000})}
\begin{itemize}
    \item \textbf{AFIO\_EXTICR1 a 4} (Offsets \texttt{0x08} a \texttt{0x14}) -- \textit{EXTI Configuration Regs}
    \begin{itemize}
        \item 4 bits por pin para elegir el puerto (0000=PA, 0001=PB, 0010=PC). \textit{Ej: EXTICR1 tiene pines 0 a 3.}
    \end{itemize}
\end{itemize}

\section*{EXTI (\texttt{0x4001 0400})}
\begin{itemize}
    \item \textbf{EXTI\_IMR} (Offset \texttt{0x00}) -- \textit{Interrupt Mask Register}
    \begin{itemize}
        \item \texttt{[19:0]} \texttt{MRy}: 1 = Habilita (desenmascara) int. $y$.
    \end{itemize}
    \item \textbf{EXTI\_RTSR / FTSR} (Offsets \texttt{0x08} / \texttt{0x0C}) -- \textit{Rising / Falling Trigger}
    \begin{itemize}
        \item \texttt{[19:0]} \texttt{TRy}: 1 = Flanco ascendente / descendente.
    \end{itemize}
    \item \textbf{EXTI\_PR} (Offset \texttt{0x14}) -- \textit{Pending Register}
    \begin{itemize}
        \item \texttt{[19:0]} \texttt{PRy}: 1 = Int. ocurrida (Limpiar con 1).
    \end{itemize}
\end{itemize}

\section*{NVIC (\texttt{0xE000 E100})}
\begin{itemize}
    \item \textbf{ISER} (\texttt{0x000}) / \textbf{ICER} (\texttt{0x080}) -- \textit{Set / Clear Enable Regs}
    \item \textbf{ISPR} (\texttt{0x100}) / \textbf{ICPR} (\texttt{0x180}) -- \textit{Set / Clear Pending Regs}
    \item \textbf{NVIC\_ISERx}:
    \begin{itemize}
        \item ISER0 (Offset \texttt{0x00}): Habilita IRQs 0 a 31 ($1 \ll \mathtt{IRQn}$).
        \item ISER1 (Offset \texttt{0x04}): Habilita IRQs 32 a 63 ($1 \ll (\mathtt{IRQn}-32)$).
    \end{itemize}
    \item \textbf{Números de Interrupción (\texttt{IRQn}):}
    \begin{itemize}
        \item \texttt{EXTI0}..\texttt{4}: 6..10 $|$ \texttt{ADC1\_2}: 18 $|$ \texttt{EXTI9\_5}: 23 $|$ \texttt{TIM2}: 28
        \item \texttt{USART1}: 37 \textit{(ISER1)} $|$ \texttt{EXTI15\_10}: 40 \textit{(ISER1)}
    \end{itemize}
\end{itemize}

\section*{SysTick (\texttt{0xE000 E010})}
\begin{itemize}
    \item \textbf{CTRL} (Offset \texttt{0x00}) -- \textit{SysTick Control and Status Register}
    \begin{itemize}
        \item \texttt{[16]} \texttt{COUNTFLAG}, \texttt{[2]} \texttt{CLKSOURCE}, \texttt{[1]} \texttt{TICKINT}, \texttt{[0]} \texttt{ENABLE}
    \end{itemize}
    \item \textbf{LOAD} (\texttt{0x04}) / \textbf{VAL} (\texttt{0x08}) / \textbf{CALIB} (\texttt{0x0C}) -- \textit{Reload / Current / Calib}
    \begin{itemize}
        \item \texttt{[23:0]}: Valor de recarga / Valor actual / Calibración.
    \end{itemize}
\end{itemize}

\section*{ADC1 (\texttt{0x4001 2400})}
\begin{itemize}
    \item \textbf{CR2} (Offset \texttt{0x08}) -- \textit{Control Register 2}
    \begin{itemize}
        \item \texttt{[23]} \texttt{TSVREFE}: Temp. \texttt{[22]} \texttt{SWSTART}: Inicia manual.
        \item \texttt{[20]} \texttt{EXTTRIG}: Disparo ext. \texttt{[19:17]} \texttt{EXTSEL}: Evento ext.
        \item \texttt{[1]} \texttt{CONT}: 1 = Continuo. \texttt{[0]} \texttt{ADON}: Enciende/Dispara.
    \end{itemize}
    \item \textbf{SR} (Offset \texttt{0x00}) -- \textit{Status Register}
    \begin{itemize}
        \item \texttt{[1]} \texttt{EOC}: 1 = Fin conversión (limpia al leer DR).
    \end{itemize}
    \item \textbf{SQR1 / SQR3} (Offsets \texttt{0x2C} / \texttt{0x34}) -- \textit{Sequence Registers}
    \begin{itemize}
        \item \texttt{SQR1[23:20]} \texttt{L}: Cantidad-1. \texttt{SQR3[4:0]} \texttt{SQ1}: 1er canal.
    \end{itemize}
    \item \textbf{DR} (Offset \texttt{0x4C}) -- \textit{Regular Data Register}
\end{itemize}

\section*{USART1 (\texttt{0x4001 3800})}
\begin{itemize}
    \item \textbf{CR1} (Offset \texttt{0x0C}) -- \textit{Control Register 1}
    \begin{itemize}
        \item \texttt{[13]} \texttt{UE}: Habilita. \texttt{[5]} \texttt{RXNEIE}: Int RX. \texttt{[3]} \texttt{TE} / \texttt{[2]} \texttt{RE}: TX/RX.
    \end{itemize}
    \item \textbf{BRR} (Offset \texttt{0x08}) -- \textit{Baud Rate Register} 
    \begin{itemize}
        \item \texttt{[15:4]} \texttt{DIV\_Mantissa} / \texttt{[3:0]} \texttt{DIV\_Fraction}.
    \end{itemize}
    \item \textbf{SR} (Offset \texttt{0x00}) -- \textit{Status Register}
    \begin{itemize}
        \item \texttt{[7]} \texttt{TXE}: Buffer vacío. \texttt{[6]} \texttt{TC}: Transmisión compl. \texttt{[5]} \texttt{RXNE}: Dato listo.
    \end{itemize}
    \item \textbf{DR} (Offset \texttt{0x04}) -- \textit{Data Register}
\end{itemize}

\section*{TIM2 - PWM (\texttt{0x4000 0000})}
\begin{itemize}
    \item \textbf{CR1} (Offset \texttt{0x00}) -- \textit{Control Register 1} (\texttt{[0]} \texttt{CEN}: Enciende)
    \item \textbf{PSC / ARR} (Offsets \texttt{0x28} / \texttt{0x2C}) -- \textit{Prescaler / Auto-Reload}
    \item \textbf{CCMR1 / CCMR2} (Offsets \texttt{0x18} / \texttt{0x1C}) -- \textit{Capture/Compare Mode}
    \begin{itemize}
        \item \texttt{[6:4]/[14:12] OCxM}: Modo (110=PWM). \texttt{[3]/[11] OCxPE}: Preload.
    \end{itemize}
    \item \textbf{CCER} (Offset \texttt{0x20}) -- \textit{Capture/Compare Enable Register}
    \begin{itemize}
        \item \texttt{[1]} \texttt{CC1P}: Salida invertida. \texttt{[0]} \texttt{CC1E}: Habilita pin.
    \end{itemize}
    \item \textbf{CCRx} (Offsets \texttt{0x34}..\texttt{0x40}) -- \textit{Capture/Compare Register}
\end{itemize}

\end{multicols*}
\end{document}
